\subsection{One-sided detection and measured footprint calibration}
The preceding stationary upper bound is not sharp as a bound for the
available commands. The fine phase need not estimate a small difference
of pooled means. Short-line randomization of nominal centers gives a
pooled forward response with a geometric zero side. Its positive mean
at depth $e$ is at least $c e^{\gamma+2}$. Testing whether a hit occurs
uses the inverse of this probability, rather than its inverse square.
Theorem~\ref{thm:line-reconstruction} consequently replaces the sufficient
power $((2\gamma+3)s+1)/(s-2)$ by
\[
                         \frac{(\gamma+2)s+1}{s-2}.
\]
This is a different command design with more spatial centers, not a
claim that the old occupation estimate has a smaller variance.

\begin{theorem}[One reference suffices for an unknown footprint]
\label{thm:reference-main}
Let the original table belong to $\mathcal B_\eta$, and let its unknown
convex launch footprint have the uniform bounds of
Section~\ref{sec:reference-calibration}. With one calibrated isolated
reference disk in the same laboratory coordinate frame, pooled collision
bits from the two scenes determine both the footprint and the table.
A finite experiment independent of the nuisance densities gives $C^2$
error $C\nu$ and correct primitive discrete data with probability at
least $1-\delta$, using at most
\[
 C\nu^{-((\gamma+2)s+1)/(s-2)}
                  \log(C/\nu)\log(C/(\nu\delta))
\]
attempts in total. The two launch densities may differ, provided they
have the same footprint and each is stationary within its scene with
the specified common boundary-mass lower bound. The footprint, its
placement and the densities are not supplied as functions. The reference
shape and pose, class bounds, nominal positioning precision, exact
compass flights and nonperiod margin remain inputs. The launch spread
is fixed as $\nu$ decreases.
\end{theorem}

The exact identification, stability modulus and finite construction
are proved in Theorems~\ref{thm:reference-exact} and
\ref{thm:reference-finite}. The reference is a disclosed additional
calibration experiment; no reference obstacle in the unknown table is
assumed known. All original localized and stationary theorems remain
active. The new upper bound is not a matching minimax characterization
of the stationary experiment.
